\documentclass[11pt,a4paper]{article}

\usepackage[a4paper,margin=2.35cm]{geometry}
\usepackage{fontspec}
\usepackage{polyglossia}
\setmainlanguage{turkish}
\setotherlanguage{english}
\setmainfont{DejaVu Serif}
\setsansfont{DejaVu Sans}
\setmonofont{DejaVu Sans Mono}
\usepackage{xcolor}
\usepackage{hyperref}
\usepackage{enumitem}
\usepackage{booktabs}
\usepackage{array}
\usepackage{longtable}
\usepackage{fancyhdr}
\usepackage{titlesec}
\usepackage{listings}
\usepackage{caption}
\usepackage{microtype}

\hypersetup{
  colorlinks=true,
  linkcolor=blue!50!black,
  urlcolor=blue!50!black,
  citecolor=blue!50!black,
  pdftitle={Task 1 SRAM PUF puf\_module.v Aciklamasi},
  pdfauthor={Alp Bolukbasi - calisma notu}
}

\definecolor{codebg}{RGB}{248,248,248}
\definecolor{codeframe}{RGB}{220,220,220}
\definecolor{keyword}{RGB}{0,70,140}
\definecolor{comment}{RGB}{90,110,90}
\definecolor{string}{RGB}{130,60,30}

\lstdefinestyle{verilogstyle}{
  basicstyle=\ttfamily\small,
  backgroundcolor=\color{codebg},
  frame=single,
  rulecolor=\color{codeframe},
  breaklines=true,
  columns=fullflexible,
  keepspaces=true,
  showstringspaces=false,
  keywordstyle=\color{keyword}\bfseries,
  commentstyle=\color{comment},
  stringstyle=\color{string},
  morekeywords={module,input,output,wire,reg,always,begin,end,if,else,case,endcase,localparam,parameter,assign,posedge,negedge,default,function,endfunction,integer,initial},
  morecomment=[l]{//},
  morecomment=[s]{/*}{*/}
}

\lstdefinestyle{bashstyle}{
  basicstyle=\ttfamily\small,
  backgroundcolor=\color{codebg},
  frame=single,
  rulecolor=\color{codeframe},
  breaklines=true,
  columns=fullflexible,
  keepspaces=true,
  showstringspaces=false
}

\titleformat{\section}{\Large\bfseries\sffamily}{\thesection}{0.6em}{}
\titleformat{\subsection}{\large\bfseries\sffamily}{\thesubsection}{0.6em}{}
\titleformat{\subsubsection}{\normalsize\bfseries\sffamily}{\thesubsubsection}{0.6em}{}
\setlist[itemize]{topsep=0.25em,itemsep=0.15em}
\setlist[enumerate]{topsep=0.25em,itemsep=0.15em}

\setlength{\headheight}{14pt}
\pagestyle{fancy}
\fancyhf{}
\lhead{Task 1 - SRAM PUF}
\rhead{puf\_module.v aciklamasi}
\cfoot{\thepage}

\title{\textbf{Task 1 SRAM PUF Tasarimi: \texttt{puf\_module.v} Teknik Aciklamasi}}
\author{Calisma ve savunma notu}
\date{\today}

\begin{document}
\maketitle

\begin{abstract}
Bu dokuman, Hardware Security Task 1 kapsaminda gelistirilen \texttt{puf\_module.v} cozumunu, bir TA'ya anlatir gibi ama teknik detaylari kaybetmeden aciklar. Odak noktasi, SRAM/BRAM tabanli PUF verisinin \texttt{combined\_ram} uzerinden okunmasi, 16-bit RAM word'lerinin 8-bit UART byte'larina bolunmesi, FSM ile bu aktarimin kontrol edilmesi ve UART handshake'inin dogru kurulmasidir. Dokuman ayrica debug amaciyla eklenen writeback ve diagnostic mode fikirlerini final PUF davranisindan ayirir.
\end{abstract}

\tableofcontents
\newpage

\section{Genel Tasarim ve Moduller Arasi Iliski}

\subsection{Buyuk resim}
\texttt{puf\_module}, PC ile FPGA icindeki RAM altyapisi arasinda duran kucuk bir kontrolcudur. Modulu tek basina bir PUF hesaplayicisi gibi dusunmek dogru degildir. Asil gorevi, FPGA icindeki SRAM/BRAM power-up icerigini bozmadan okumaya calismak ve bunu UART uzerinden PC'ye byte stream olarak gondermektir.

\begin{center}
\begin{tabular}{c c c c c}
PC Python script & $\longleftrightarrow$ & UART & $\longleftrightarrow$ & \texttt{puf\_module FSM} $\longleftrightarrow$ \texttt{combined\_ram} \\
\end{tabular}
\end{center}

Normal akis soyledir:
\begin{enumerate}
  \item PC, UART uzerinden \texttt{s} veya \texttt{S} komutunu gonderir.
  \item \texttt{puf\_module}, UART RX tarafindan bu komutu gorur.
  \item FSM, \texttt{combined\_ram} adres alanini bastan sona okur.
  \item RAM 16-bit word dondurdugu icin her word iki adet 8-bit UART byte'ina bolunur.
  \item UART TX bu byte'lari PC'ye yollar.
  \item Python script toplam 16384 byte'i 512 satir, satir basina 64 hex karakter olarak dosyaya yazar.
\end{enumerate}

Bu tasarimda final hedef sudur:
\begin{lstlisting}[style=bashstyle]
512 PUF instance * 256 bit = 131072 bit = 16384 byte = 16 KiB
\end{lstlisting}

\subsection{Portlar: dis dunya ile konusulan sinyaller}
\texttt{puf\_module} portlari iki gruba ayrilabilir: clock/reset ve UART arayuzu.

\begin{lstlisting}[style=verilogstyle]
input  clk,
input  rst,
input  uart_rx_ready,
input  [7:0] uart_data_from_rx,
input  uart_tx_ready,
output [7:0] uart_data_to_tx,
output uart_tx_enable
\end{lstlisting}

\begin{longtable}{>{\ttfamily}p{4.2cm}p{10.4cm}}
\toprule
Sinyal & Anlam \\
\midrule
\endhead
clk & Sequential FSM'in ilerledigi ana clock. \\
rst & FSM ve ic register'lari bilinen duruma donduren reset. RAM icerigini sifirlamaz. \\
uart\_rx\_ready & PC'den yeni bir UART byte'i geldigini bildirir. \\
uart\_data\_from\_rx & RX tarafindan yakalanmis komut byte'i. \texttt{s/S} normal PUF dump, \texttt{d/D} debug dump. \\
uart\_tx\_ready & UART TX'in yeni byte kabul etmeye hazir oldugunu bildirir. \\
uart\_data\_to\_tx & UART'a verilecek 8-bit byte. Genelde \texttt{puf\_byte\_reg}'den gelir. \\
uart\_tx\_enable & UART'a ``su byte'i al ve gonder'' diyen bir-cycle pulse. \\
\bottomrule
\end{longtable}

\subsection{RAM arayuzu ve \texttt{combined\_ram}}
\texttt{puf\_module}, fiziksel RAM bloklarini dogrudan yonetmez. Bunun yerine saglanan \texttt{combined\_ram} wrapper'ina baglanir:

\begin{lstlisting}[style=verilogstyle]
combined_ram ram_inst (
   .clk   ( clk ),
   .rdata ( rdata ),
   .raddr ( raddr ),
   .we    ( we ),
   .wdata ( wdata ),
   .wmask ( wmask ),
   .waddr ( waddr )
);
\end{lstlisting}

Bu wrapper, cok sayida fiziksel \texttt{SB\_RAM40\_4K} RAM blogunu tek buyuk adres alani gibi gostermeyi amaclar. \texttt{puf\_module} acisindan interface soyledir:

\begin{lstlisting}[style=verilogstyle]
wire [15:0] rdata;
reg  [12:0] raddr;
reg  [12:0] waddr;
reg  [15:0] wmask;
reg         we;
reg  [15:0] wdata;
\end{lstlisting}

\texttt{rdata}, \texttt{combined\_ram}'den gelen 16-bit output'tur. \texttt{raddr} 13-bit read address'tir. 13-bit adres alani su anlama gelir:

\begin{lstlisting}[style=bashstyle]
2^13 = 8192 word
8192 word * 16 bit = 131072 bit = 16 KiB
\end{lstlisting}

\texttt{combined\_ram} icinde adres iki parcaya ayrilir:
\begin{lstlisting}[style=verilogstyle]
raddr[12:8]  // hangi RAM block?
raddr[7:0]   // o block icindeki local word offset
\end{lstlisting}

\texttt{puf\_module} seviyesinde ise bu detay soyutlanmistir. Biz sadece \texttt{raddr=0..8191} arasinda gezeriz ve her adreste 16-bit \texttt{rdata} bekleriz.

\subsection{Byte index mental modeli: \texttt{i\_r}, \texttt{raddr}, \texttt{rdata}}
RAM 16-bit word dondurur, UART ise 8-bit byte gonderir. Bu yuzden tek bir sayaç hem RAM adresini hem de byte secimini belirler:

\begin{lstlisting}[style=verilogstyle]
raddr = i_r[13:1];
\end{lstlisting}

\begin{center}
\begin{tabular}{ccc}
\toprule
\texttt{i\_r} byte index & \texttt{raddr} word address & Gonderilecek byte \\
\midrule
0 & 0 & \texttt{rdata[7:0]} \\
1 & 0 & \texttt{rdata[15:8]} \\
2 & 1 & \texttt{rdata[7:0]} \\
3 & 1 & \texttt{rdata[15:8]} \\
\bottomrule
\end{tabular}
\end{center}

Buradaki en onemli cumle sudur:

\begin{quote}
\texttt{i\_r} UART byte index'idir. \texttt{i\_r[13:1]} RAM word adresi olur, \texttt{i\_r[0]} ise 16-bit word icindeki lower veya upper byte secimini yapar.
\end{quote}

\subsection{State tanimlari ve \texttt{localparam}}
FSM state'leri \texttt{localparam} ile isimlendirilir:

\begin{lstlisting}[style=verilogstyle]
localparam
   INIT = 0,
   WAIT_FOR_REQUEST = 1,
   WAITCYCLE_FOR_MEMORY = 2,
   PUF_READ = 3,
   UART_SEND = 4,
   UART_WAIT_FINISH = 5,
   LOOP_CONDITION = 6;

reg [2:0] state;
\end{lstlisting}

7 state oldugu icin 3-bit \texttt{state} yeterlidir. \texttt{localparam} kullanmak, \texttt{state <= 3;} gibi okunmasi zor kod yerine \texttt{state <= PUF\_READ;} yazmayi saglar.

\begin{longtable}{>{\ttfamily}p{4.7cm}p{9.9cm}}
\toprule
State & Gorev \\
\midrule
\endhead
INIT & Reset sonrasi baslangic ve dummy UART davranisi. \\
WAIT\_FOR\_REQUEST & PC'den \texttt{s/S} veya \texttt{d/D} bekler. \\
WAITCYCLE\_FOR\_MEMORY & RAM adresinin oturmasi / synchronous RAM latency icin ara state. \\
PUF\_READ & \texttt{rdata}'dan dogru byte'i secip \texttt{puf\_byte\_reg}'e yazar. \\
UART\_SEND & UART'a bir cycle \texttt{uart\_tx\_enable} pulse verir. \\
UART\_WAIT\_FINISH & UART byte gonderimini bitirene kadar bekler. \\
LOOP\_CONDITION & Sayaç artirir veya transfer bittiyse idle state'e doner. \\
\bottomrule
\end{longtable}

\subsection{Register'lar ve initialization}
FSM'in hafizasi su register'lardan olusur:

\begin{lstlisting}[style=verilogstyle]
reg [13:0] i_r;
reg [12:0] i_w_r;
reg [7:0]  puf_byte_reg;
reg        debug_mode;
\end{lstlisting}

Reset sirasinda:
\begin{lstlisting}[style=verilogstyle]
state <= INIT;
i_r <= 14'd0;
i_w_r <= 13'd0;
debug_mode <= 1'b0;
puf_byte_reg <= 8'd0;
\end{lstlisting}

Bu initialization'lar RAM'i sifirlamak icin degildir. Sadece FSM ve kontrol register'larini bilinen duruma getirir. PUF mantigi acisindan kritik olan nokta sudur: RAM power-up icerigi korunmali ve okunmalidir; reset sadece kontrolcu tarafini temizlemelidir.

\section{Sequential Logic: Clock Edge'de Ne Oluyor?}

Sequential block su sekildedir:
\begin{lstlisting}[style=verilogstyle]
always @ (posedge clk, posedge rst) begin
   ...
end
\end{lstlisting}

Bu blok register update'lerini yapar. \texttt{state}, \texttt{i\_r}, \texttt{i\_w\_r}, \texttt{debug\_mode}, \texttt{puf\_byte\_reg} gibi degerler clock edge'lerinde guncellenir. Bu yuzden sequential block icinde genelde non-blocking assignment, yani \texttt{<=}, kullanilir.

\subsection{Reset davranisi}
Reset geldiginde FSM ve sayaçlar bilinen duruma doner:
\begin{lstlisting}[style=verilogstyle]
if (rst) begin
   state <= INIT;
   i_r <= 14'd0;
   i_w_r <= 13'd0;
   debug_mode <= 1'b0;
   puf_byte_reg <= 8'd0;
end
\end{lstlisting}

Bu davranis TA'ya soyle aciklanabilir:

\begin{quote}
Reset, PUF memory'yi initialize etmez. Sadece transfer FSM'ini ve internal control register'larini initialize eder. Boylece yeni request geldiginde capture bastan ve deterministic bir kontrol akisiyle baslar.
\end{quote}

\subsection{\texttt{INIT}}
\texttt{INIT}, bir baslangic state'idir:
\begin{lstlisting}[style=verilogstyle]
INIT : begin
   state <= WAIT_FOR_REQUEST;
end
\end{lstlisting}

Sequential tarafta tek gorevi FSM'i request bekleme durumuna almaktir. Combinational tarafta ise dummy UART byte gonderme davranisi olabilir.

\subsection{\texttt{WAIT\_FOR\_REQUEST}}
Bu state idle durumudur. FSM burada PC'den komut bekler.

\begin{lstlisting}[style=verilogstyle]
i_r <= 14'd0;

if (uart_rx_ready &&
   (uart_data_from_rx == "s" || uart_data_from_rx == "S")) begin
   debug_mode <= 1'b0;
   state <= WAITCYCLE_FOR_MEMORY;
end else if (uart_rx_ready &&
   (uart_data_from_rx == "d" || uart_data_from_rx == "D")) begin
   debug_mode <= 1'b1;
   state <= WAITCYCLE_FOR_MEMORY;
end
\end{lstlisting}

Burada \texttt{i\_r}'nin sifirlanmasi onemlidir. Her yeni capture RAM adres alaninin basindan baslamalidir. \texttt{s/S} normal PUF dump'i baslatir. \texttt{d/D} ise debug packet modunu baslatir.

Debug icin eklenen \texttt{i\_w\_r} artirma davranisi, final PUF fikrinin parcasi degildir. Sadece RAM write/readback yolunu test etmeye yarayan gecici bir tekniktir.

\subsection{\texttt{WAITCYCLE\_FOR\_MEMORY}}
Bu state timing amaclidir:
\begin{lstlisting}[style=verilogstyle]
WAITCYCLE_FOR_MEMORY : begin
   state <= PUF_READ;
end
\end{lstlisting}

\texttt{raddr}, combinational logic tarafinda \texttt{i\_r[13:1]} olarak surulur. RAM synchronous oldugu icin, adresi verip ayni logical anda \texttt{rdata}'yi sample etmek risklidir. Bu state, RAM output'unun hazirlanmasi icin bir ara cycle saglar.

\subsection{\texttt{PUF\_READ}}
Bu state'te RAM'den gelen 16-bit \texttt{rdata}'dan bir byte secilir.

\begin{lstlisting}[style=verilogstyle]
if (i_r[0] == 1'b0) begin
   puf_byte_reg <= rdata[7:0];
end else begin
   puf_byte_reg <= rdata[15:8];
end
\end{lstlisting}

Secilen byte hemen UART'a verilmez. Once \texttt{puf\_byte\_reg} register'ina alinir. Bir sonraki state olan \texttt{UART\_SEND}, bu stable byte'i UART'a verir.

Debug mode'da ise \texttt{rdata} yerine internal signal snapshot'lari ve marker byte'lari gonderilir:
\begin{lstlisting}[style=verilogstyle]
6'd0:  puf_byte_reg <= 8'h44; // 'D'
6'd1:  puf_byte_reg <= 8'h42; // 'B'
...
6'd9:  puf_byte_reg <= rdata[7:0];
6'd10: puf_byte_reg <= rdata[15:8];
...
\end{lstlisting}

Bu sayede UART/FSM path'i ile RAM readout path'i birbirinden ayrilabilir.

\subsection{\texttt{UART\_SEND}}
Sequential block'ta:
\begin{lstlisting}[style=verilogstyle]
UART_SEND : begin
   state <= UART_WAIT_FINISH;
end
\end{lstlisting}

Bu state bir cycle suren pulse state'tir. Asil \texttt{uart\_tx\_enable} combinational block'ta bu state sirasinda 1 yapilir. Sequential block hemen \texttt{UART\_WAIT\_FINISH}'e gecisi planlar.

\subsection{\texttt{UART\_WAIT\_FINISH}}
\begin{lstlisting}[style=verilogstyle]
UART_WAIT_FINISH : begin
   if (uart_tx_ready) begin
      state <= LOOP_CONDITION;
   end
end
\end{lstlisting}

UART fiziksel olarak yavas oldugu icin bir byte'in gonderimi bircok FPGA clock cycle'i surebilir. Bu state, UART'in tekrar hazir olmasini bekler. Bu olmazsa FSM byte'lari UART'in kabul edebileceginden hizli gonderir.

\subsection{\texttt{LOOP\_CONDITION}}
Bu state transferin devam edip etmeyecegine karar verir:
\begin{lstlisting}[style=verilogstyle]
if ((debug_mode && i_r == DEBUG_BYTES - 1) ||
    (!debug_mode && i_r == PUF_BYTES - 1)) begin
   debug_mode <= 1'b0;
   state <= WAIT_FOR_REQUEST;
end else begin
   i_r <= i_r + 1'b1;
   state <= WAITCYCLE_FOR_MEMORY;
end
\end{lstlisting}

Normal mode'da \texttt{i\_r} 0'dan 16383'e gider. Debug mode'da daha kisa bir diagnostic packet gonderilir. Son byte gonderilince FSM tekrar \texttt{WAIT\_FOR\_REQUEST}'e doner.

\section{Combinational Logic: Output ve Control Sinyalleri}

Combinational block su sekildedir:
\begin{lstlisting}[style=verilogstyle]
always @ (*) begin
   ...
end
\end{lstlisting}

Bu blok clock beklemez. O anki \texttt{state}, counter'lar ve input'lara gore output/control sinyallerini hazirlar.

\subsection{Default assignment felsefesi}
Combinational block'un basinda safe default degerler verilir:
\begin{lstlisting}[style=verilogstyle]
uart_tx_enable <= 'b0;
uart_data_to_tx = puf_byte_reg;
raddr = i_r[13:1];
wmask <= 0;
waddr = 13'd0;
wdata = 16'd0;
we = 1'b0;
\end{lstlisting}

Bu default'lar latch cikmasini engeller ve sinyalleri guvenli durumda tutar. Mesela \texttt{we} default olarak 0 olmalidir; aksi halde RAM icerigi yanlislikla overwrite edilebilir.

\subsection{UART output sinyalleri}
\texttt{uart\_data\_to\_tx} normalde \texttt{puf\_byte\_reg}'e baglidir:
\begin{lstlisting}[style=verilogstyle]
uart_data_to_tx = puf_byte_reg;
\end{lstlisting}

\texttt{uart\_tx\_enable} default olarak 0'dir. Sadece \texttt{UART\_SEND} state'inde 1 yapilir:
\begin{lstlisting}[style=verilogstyle]
UART_SEND : begin
   uart_tx_enable = 1'b1;
end
\end{lstlisting}

Bu, UART icin bir-cycle transmit pulse anlamina gelir.

\subsection{RAM read address}
\begin{lstlisting}[style=verilogstyle]
raddr = i_r[13:1];
\end{lstlisting}

Bu assignment surekli aktiftir. FSM \texttt{i\_r}'yi her byte sonrasinda artirdikca, \texttt{raddr} otomatik olarak dogru RAM word adresini gosterir. \texttt{i\_r[0]} ise lower/upper byte secimi icin sequential block'ta kullanilir.

\subsection{RAM write sinyalleri ve debug writeback}
Normal PUF readout read-only olmalidir. Bu nedenle default:
\begin{lstlisting}[style=verilogstyle]
we = 1'b0;
waddr = 13'd0;
wdata = 16'd0;
wmask <= 0;
\end{lstlisting}

Debug writeback modunda, \texttt{WAIT\_FOR\_REQUEST} sirasinda RAM'e gorunur bir pattern yazmayi denedik:
\begin{lstlisting}[style=verilogstyle]
waddr = i_w_r;
wdata = 16'hA5A5;
we = 1'b1;
\end{lstlisting}

Bunun amaci final cozum degil, sadece \texttt{combined\_ram} write/read path'inin calisip calismadigini test etmektir.

\subsection{State-by-state combinational behavior}
\begin{itemize}
  \item \texttt{INIT}: dummy byte gondermek icin \texttt{uart\_tx\_enable=1}, \texttt{uart\_data\_to\_tx=0}.
  \item \texttt{WAIT\_FOR\_REQUEST}: normalde default'lar. Debug writeback aciksa \texttt{waddr/wdata/we} surulur.
  \item \texttt{WAITCYCLE\_FOR\_MEMORY}: default'lar yeterlidir; \texttt{raddr} zaten \texttt{i\_r}'ye baglidir.
  \item \texttt{PUF\_READ}: asıl byte sampling sequential block'ta yapilir; combinational tarafta default'lar yeterlidir.
  \item \texttt{UART\_SEND}: \texttt{uart\_tx\_enable=1} pulse verilir.
  \item \texttt{UART\_WAIT\_FINISH}: default'lar; enable tekrar 0 olur.
  \item \texttt{LOOP\_CONDITION}: karar sequential block'ta verilir; combinational tarafta default'lar.
\end{itemize}

\section{PUF-UART Handshake ve Tam Transfer Akisi}

\subsection{Request handshake}
Python script UART uzerinden komut byte'i yollar. FPGA tarafinda \texttt{uart\_rx\_ready} 1 oldugunda \texttt{uart\_data\_from\_rx} gecerli kabul edilir. FSM sadece \texttt{WAIT\_FOR\_REQUEST} state'inde bu komutu dikkate alir.

\begin{lstlisting}[style=verilogstyle]
if (uart_rx_ready && uart_data_from_rx == "s") begin
   debug_mode <= 1'b0;
   state <= WAITCYCLE_FOR_MEMORY;
end
\end{lstlisting}

Burada request, tum transferi baslatan tetiktir.

\subsection{RAM'den byte hazirlama}
Request geldikten sonra \texttt{i\_r=0} iken \texttt{raddr=0} olur. \texttt{WAITCYCLE\_FOR\_MEMORY} sonrasinda \texttt{PUF\_READ} state'i \texttt{rdata}'nin lower veya upper byte'ini \texttt{puf\_byte\_reg}'e alir.

\begin{lstlisting}[style=verilogstyle]
RAM rdata[15:0] -> puf_byte_reg[7:0] -> uart_data_to_tx[7:0]
\end{lstlisting}

\texttt{puf\_byte\_reg}, RAM ile UART arasinda stable bir byte buffer'i gibi davranir.

\subsection{UART TX handshake}
UART TX modeli soyledir:
\begin{enumerate}
  \item \texttt{uart\_tx\_ready=1} ise UART yeni byte kabul edebilir.
  \item \texttt{uart\_data\_to\_tx} stable tutulur.
  \item \texttt{uart\_tx\_enable} bir cycle 1 yapilir.
  \item UART byte'i alir ve fiziksel olarak serial hatta gondermeye baslar.
  \item FSM \texttt{UART\_WAIT\_FINISH}'te \texttt{uart\_tx\_ready}'nin tekrar 1 olmasini bekler.
\end{enumerate}

Bu handshake olmadan FSM, UART'in kaldirabileceginden daha hizli byte gonderebilir.

\subsection{Bir byte icin timing}
Bir byte'in tam akisi soyledir:

\begin{longtable}{p{3cm}p{11.4cm}}
\toprule
Asama & Olay \\
\midrule
\endhead
\texttt{WAITCYCLE\_FOR\_MEMORY} & \texttt{raddr = i\_r[13:1]} RAM tarafinda stable hale gelir. \\
\texttt{PUF\_READ} & \texttt{rdata}'dan lower/upper byte secilip \texttt{puf\_byte\_reg}'e yazilir. \\
\texttt{UART\_SEND} & \texttt{uart\_data\_to\_tx = puf\_byte\_reg}; \texttt{uart\_tx\_enable=1}. \\
\texttt{UART\_WAIT\_FINISH} & UART busy iken beklenir; \texttt{uart\_tx\_ready} tekrar 1 olunca devam edilir. \\
\texttt{LOOP\_CONDITION} & \texttt{i\_r} artirilir veya transfer bittiyse idle'a donulur. \\
\bottomrule
\end{longtable}

\subsection{Debug packet ve stream alignment problemi}
Debug icin \texttt{d/D} komutuyla marker byte'lari gonderildi:
\begin{lstlisting}[style=bashstyle]
44 42 = 'D' 'B'
A5 5A C3 3C = bilinen marker pattern
\end{lstlisting}

Bu marker'larin amaci, UART/FSM path'inin calistigini ve stream icindeki hizalanma problemlerini gormektir. 16 KiB debug capture'da once uzun \texttt{00} bolgesi, sonra \texttt{44 42}, sonra \texttt{EE} default byte'lari gorulmesi su hipotezi dogurdu:

\begin{itemize}
  \item Python read window'u FPGA stream baslangicina tam hizali olmayabilir.
  \item Reset/dummy/stale byte'lar stream basinda yer kapliyor olabilir.
  \item Debug packet geliyor, fakat stream icinde offsetli geliyor olabilir.
\end{itemize}

Normal PUF capture'da ``ilk nonzero'dan basla'' gibi bir strateji dogru degildir; cunku gercek PUF verisi \texttt{00} ile baslayabilir. Ancak debug modunda bilincli \texttt{44 42} marker'i oldugu icin marker aramak ve packet'i onun etrafinda analiz etmek mantiklidir.

\section{Sonuc ve TA'ya Anlatim Ozeti}

Bu cozumde \texttt{puf\_module}, UART komutuyla tetiklenen bir FSM olarak tasarlanmistir. \texttt{s/S} komutu normal PUF dump'i, \texttt{d/D} komutu diagnostic dump'i baslatir. RAM 16-bit word dondurdugu icin \texttt{i\_r} byte counter'i hem RAM word adresini hem de lower/upper byte secimini belirler. Secilen byte \texttt{puf\_byte\_reg}'e alinir, sonra UART TX handshake ile PC'ye gonderilir. Normal modda 16384 byte, yani 512 adet 256-bit PUF instance'i gonderilir.

Kisa savunma cumlesi:

\begin{quote}
I implemented the FSM as a bridge between the combined SRAM/BRAM interface and the UART transmitter. The RAM side is word-oriented and returns 16-bit data, while the UART side is byte-oriented and requires a ready/enable handshake. Therefore, I use a byte counter \texttt{i\_r}; \texttt{i\_r[13:1]} addresses the RAM word and \texttt{i\_r[0]} selects the low or high byte. The selected byte is registered in \texttt{puf\_byte\_reg}, then transmitted only when \texttt{uart\_tx\_ready} is high. The FSM waits for the UART to finish before moving to the next byte. Debug modes were added only to distinguish UART/FSM problems from RAM readout issues.
\end{quote}

\end{document}
